\documentclass[12pt]{article}
\usepackage[utf8]{inputenc}
\usepackage[T1]{fontenc}
\usepackage[margin=1in]{geometry}
\usepackage{setspace}
\usepackage{xcolor}
\usepackage[colorlinks=true, urlcolor=blue!60!black, linkcolor=blue!60!black]{hyperref}
\setlength{\parindent}{2em}
\setlength{\parskip}{1.25em}
\renewcommand{\baselinestretch}{1.0}
\begin{document}
\begin{center}
{\LARGE \textit{San Antonio Independent School District v.\ Rodriguez}, 411 U.S.\ 1 (1973)\par}
\vspace{0.8em}
{\large Case Brief --- Four-Page Edition\par}
\end{center}
\vspace{1em}

\section{Memory Jogger}

Texas's property-tax-based school funding did not violate equal protection: neither poverty nor district wealth is a suspect classification, and education is not a fundamental federal constitutional right (\href{https://supreme.justia.com/cases/federal/us/411/1/}{Justia}).

\section{Detailed Case Facts}

Texas funded its schools through the state Minimum Foundation School Program, which guaranteed a basic minimum education in every district, plus local property taxes capped at \$1.50 per \$100 of valuation. Because property values varied widely, spending did too. In 1967--1968, Edgewood ISD in San Antonio, about 90\% Mexican-American with the area's lowest property value per pupil (\$5,960), spent \$356 per pupil despite the highest tax rate. Nearby Alamo Heights, affluent and predominantly ``Anglo,'' with more than \$49,000 in property per pupil, spent \$594 at a lower rate, and even received more state aid per pupil because the formula rewarded experienced teachers. Demetrio P.\ Rodriguez and other Edgewood parents sued on behalf of poor and minority schoolchildren in low-property-wealth districts throughout Texas.

\section{Procedural History}

Plaintiffs sued in 1968 in the Western District of Texas. In December 1971, a three-judge court held that wealth is a suspect classification and education a fundamental interest, and that the system failed both strict scrutiny and rational basis review; it stayed its mandate to give Texas time to legislate. On direct appeal, the Supreme Court reversed on March 21, 1973.

\section{Judicial Votes}

Reversed, \textbf{5--4}.
\begin{itemize}
\item \textbf{Majority:} Justice Powell, joined by Chief Justice Burger and Justices Stewart, Blackmun, and Rehnquist.
\item \textbf{Concurrence:} Justice Stewart.
\item \textbf{Dissents:} Justice Brennan; Justice White, joined by Douglas and Brennan; Justice Marshall, joined by Douglas.
\end{itemize}

\section{Holding}

Strict scrutiny applies only to laws disadvantaging a suspect class or impinging on a right explicitly or implicitly guaranteed by the Constitution. District-wealth disparities create no suspect class where there is no absolute deprivation and no definable class of indigents, and education is not a fundamental right. Texas's system, which guarantees a basic education while relying on local taxes for local control, is rational. The suggestion that an \emph{absolute deprivation} of education, or failure to provide ``an opportunity to acquire the basic minimal skills necessary'' for speech and voting, might be different is dicta, later tested in \textit{Plyler v.\ Doe}.

\section{Analysis of Opinions}

\textbf{Majority (Powell).} Wealth cases like \textit{Griffin} and \textit{Douglas v.\ California} involved a definable class unable to pay and an absolute deprivation. Texas guaranteed twelve years of free schooling, and plaintiffs' own data did not show that poor families lived in property-poor districts. ``[T]he importance of a service performed by the State does not determine whether it must be regarded as fundamental'': as with housing (\textit{Lindsey}) and welfare (\textit{Dandridge}), the test is constitutional text. The nexus theory proved too much, since food and shelter also matter to civic life, and the Constitution does not guarantee ``the most effective speech or the most informed electoral choice.'' The Court's lack of expertise and federalism counseled restraint, and local control rationally justified the system even with ``some inequality.''

\textbf{Stewart (concurring).} The system is ``chaotic and unjust,'' but the Equal Protection Clause ``confers no substantive rights and creates no substantive liberties.''

\textbf{Brennan (dissenting).} Education is ``inextricably linked'' to voting and speech, so strict scrutiny applies.

\textbf{White (dissenting).} The system fails even rational basis review: Edgewood would need a \$5.76 rate, far above the \$1.50 cap, to match the highest yield, so ``local control'' is a real option only for rich districts.

\textbf{Marshall (dissenting).} The decision is ``a retreat from our historic commitment to equality of educational opportunity.'' Marshall rejected the rigid two tiers for a ``spectrum of standards'' tied to the importance of the interest and its nexus to constitutional rights, and called local control a ``mere sham.''

\section{Examples: Future Applications}

\begin{enumerate}
\item \textbf{Same-Side: Bus Fees.} A modest school-bus fee without an indigency waiver, where children still reach school, gets rational basis review, as in \textit{Kadrmas} (1988).
\item \textbf{Same-Side: Libraries and Parks.} Funding them by local property tax creates disparities, but they are not constitutional rights.
\item \textbf{Opposite-Side: Excluding Undocumented Children.} Withholding all education triggers the absolute-deprivation caveat, as in \textit{Plyler v.\ Doe} (1982).
\item \textbf{Opposite-Side: Appeal Fees.} A \$500 fee barring indigent defendants from a first appeal is an absolute wealth-based deprivation under \textit{Griffin} and \textit{Douglas}.
\item \textbf{Fence-Sitter: ``Access to Literacy.''} Students in a collapsing district claim denial of basic minimal skills, as in \textit{Gary B.\ v.\ Whitmer} (6th Cir.\ 2020, vacated), invoking Powell's reservation against his broader holding.
\end{enumerate}

\section{Critique}

\textbf{Federalism.} Judge Jeffrey S.\ Sutton argues that restraint spurred state-court litigation under state education clauses, including Texas's \textit{Edgewood} cases, and others note the lack of manageable standards. Derek W.\ Black offers an originalist challenge: Reconstruction Congresses required States to guarantee education as a condition of readmission.

\textbf{Progressive.} Erwin Chemerinsky and critical race scholars say the Court entrenched ``separate and unequal'' schools by treating race and wealth as separate, and scholars note the two-tier model soon gave way to intermediate scrutiny and \textit{Plyler}'s heightened rationality.

\section{Key Quotations}

\begin{quote}
``The system of alleged discrimination and the class it defines have none of the traditional indicia of suspectness: the class is not saddled with such disabilities, or subjected to such a history of purposeful unequal treatment, or relegated to such a position of political powerlessness as to command extraordinary protection from the majoritarian political process.''
\end{quote}
\noindent --- Justice Powell, for the Court.

\begin{quote}
``It is not the province of this Court to create substantive constitutional rights in the name of guaranteeing equal protection of the laws. Thus, the key to discovering whether education is `fundamental' is not to be found in comparisons of the relative societal significance of education, as opposed to subsistence or housing. Nor is it to be found by weighing whether education is as important as the right to travel. Rather, the answer lies in assessing whether there is a right to education explicitly or implicitly guaranteed by the Constitution.''
\end{quote}
\noindent --- Justice Powell, for the Court.

\begin{quote}
``Education, of course, is not among the rights afforded explicit protection under our Federal Constitution. Nor do we find any basis for saying it is implicitly so protected.''
\end{quote}
\noindent --- Justice Powell, for the Court.

\end{document}
